\SourcePage{6}

\bdsec{1.1. 初始定义}

关于叠的一般参考文献是 [LMB93]。

\bdp{1.1.1} 设 $\mathcal{Y}$ 是 $\mathbb{C}$ 上一个光滑的等维代数叠。记 $\Theta_{\mathcal{Y}}$ 为切层；这是 $\mathcal{Y}$ 上的凝聚层。余切叠 $T^{*}\mathcal{Y}=\operatorname{Spec}\,\operatorname{Sym}\Theta_{\mathcal{Y}}$ 未必光滑。一般而言，$\dim T^{*}\mathcal{Y}=2\dim\mathcal{Y}$ 也不成立（例如考虑无限代数群的分类叠，或 $\mathfrak{sl}_n$ 模去 $\mathrm{SL}_n$ 的伴随作用所得的商）。然而总有
\begin{equation}\tag{1}
\dim T^{*}\mathcal{Y}\ge 2\dim\mathcal{Y}
\end{equation}
若
\begin{equation}\tag{2}
\dim T^{*}\mathcal{Y}=2\dim\mathcal{Y}
\end{equation}
则称 $\mathcal{Y}$ 为良。此时 $T^{*}\mathcal{Y}$ 局部是纯维数 $2\dim\mathcal{Y}$ 的完全交。若 $\mathcal{Y}=K\backslash S$，其中 $S$ 是带代数群 $K$ 作用的光滑簇，则此事显然。此时 $T^{*}\mathcal{Y}$ 由 $T^{*}S$ 经哈密顿约化得到（见 1.2.1）。一般情形与此十分相似。

容易证明，(2) 等价于以下条件：
\begin{equation}\tag{3}
\operatorname{codim}\{y\in\mathcal{Y}\mid \dim G_y=n\}\ge n\qquad\text{对一切 }n>0.
\end{equation}
这里 $G_y$ 是 $y$ 的自同构群（回忆：叠的一点可以具有非平凡对称性）。若
\begin{equation}\tag{4}
\operatorname{codim}\{y\in\mathcal{Y}\mid \dim G_y=n\}>n\qquad\text{对一切 }n>0.
\end{equation}
则称 $\mathcal{Y}$ 为极良。

容易看出，$\mathcal{Y}$ 极良当且仅当 $T^{*}\mathcal{Y}^{0}$ 在 $T^{*}\mathcal{Y}$ 中稠密；其中
\[
\mathcal{Y}^{0}:=\{y\in\mathcal{Y}\mid \dim G_y=0\}
\]
是 $\mathcal{Y}$ 的最大 Deligne--Mumford 子叠。特别地，若 $\mathcal{Y}$ 极良，则对 $\mathcal{Y}$ 的每个连通分支 $\mathcal{Y}_i$，$T^{*}\mathcal{Y}_i$ 不可约。

\emph{注.} 「良」实际上是「good for lazybones」（对懒人友好）的意思（见 1.1.4 末的注）。
% [Tutor Note: 原书 p.6 的 good / very good 是作者的戏称：「good」实指「good for lazybones」。术语表定为「良 / 极良」，此处保留原文字面以便对照。]

\SourcePage{7}

\bdp{1.1.2} 记 $\mathcal{Y}_{sm}$ 为 $\mathcal{Y}$ 的光滑拓扑（见 [LMB93, Section 6]）。$\mathcal{Y}_{sm}$ 的一个对象是光滑 1-态射 $\pi_S:S\to\mathcal{Y}$，其中 $S$ 是概形。态射 $(S,\pi_S)\to(S',\pi_{S'})$ 是一对 $(\phi,\alpha)$：$\phi:S\to S'$ 是概形的光滑态射，$\alpha$ 是 2-态射 $\pi_S\Rightarrow\pi_{S'}\phi$。我们常把 $(S,\pi_S)$ 简记为 $S$。

对 $S\in\mathcal{Y}_{sm}$，有相对切层 $\Theta_{S/\mathcal{Y}}$，它是局部自由的 $\mathcal{O}_S$-模。它位于一个典范正合列中
\[
\Theta_{S/\mathcal{Y}}\to\Theta_S\to\pi_S^{*}\Theta_{\mathcal{Y}}\to 0.
\]
因此 $\pi_S^{*}\operatorname{Sym}\Theta_{\mathcal{Y}}=\operatorname{Sym}\Theta_S/I^{cl}$，其中 $I^{cl}:=(\operatorname{Sym}\Theta_S)\Theta_{S/\mathcal{Y}}$。把 $\operatorname{Sym}\Theta_S$ 视为 $S$ 的平展拓扑上的层，它带有通常的泊松括号 $\{\,\}$。设 $\tilde{P}\subset\operatorname{Sym}\Theta_S$ 是理想 $I^{cl}$ 的 $\{\,\}$-正规化子。令 $(P_{\mathcal{Y}})_S:=\tilde{P}/I^{cl}$，于是 $(P_{\mathcal{Y}})_S$ 是 $\operatorname{Sym}\Theta_S$ 对 $\Theta_{S/\mathcal{Y}}$ 的哈密顿约化。这是 $S_{\text{ét}}$ 上分次泊松代数的层。若 $S\to S'$ 是 $\mathcal{Y}_{sm}$ 中的态射，则 $(P_{\mathcal{Y}})_S$ 等于 $(P_{\mathcal{Y}})_{S'}$ 的层论拉回。于是当 $S$ 变化时，$(P_{\mathcal{Y}})_S$ 构成 $\mathcal{Y}_{sm}$ 上泊松代数的层 $P_{\mathcal{Y}}$，称为 $\mathcal{Y}$ 的象征代数。交换代数的嵌入 $P_{\mathcal{Y}}\hookrightarrow\operatorname{Sym}\Theta_{\mathcal{Y}}$ 给出整体截面空间之间的同构
\begin{equation}\tag{5}
\Gamma(\mathcal{Y},P_{\mathcal{Y}})\simeq\Gamma(\mathcal{Y},\operatorname{Sym}\Theta_{\mathcal{Y}})=\Gamma(T^{*}\mathcal{Y},\mathcal{O}).
\end{equation}

\bdp{1.1.3} 对 $S\in\mathcal{Y}_{sm}$，考虑微分算子的层 $\mathcal{D}_S$。这是 $S_{\text{ét}}$ 上结合代数的层。设 $\tilde{D}_S\subset\mathcal{D}_S$ 是左理想 $I:=\mathcal{D}_S\Theta_{S/\mathcal{Y}}\subset\mathcal{D}_S$ 的正规化子。令 $(D_{\mathcal{Y}})_S:=\tilde{D}_S/I$。此代数从右作用在 $\mathcal{D}_S$-模 $(\mathcal{D}_{\mathcal{Y}})_S:=\mathcal{D}_S/I$ 上；该作用把 $(D_{\mathcal{Y}})_S$ 与 $\operatorname{End}_{\mathcal{D}_S}((\mathcal{D}_{\mathcal{Y}})_S)$ 的反代数等同起来。

对 $\mathcal{Y}_{sm}$ 中的任意态射 $(\phi,\alpha):S\to S'$，有显然的 $\mathcal{D}_S$-模同构 $\phi^{*}((\mathcal{D}_{\mathcal{Y}})_{S'})\simeq(\mathcal{D}_{\mathcal{Y}})_S$。它把 $(\mathcal{D}_{\mathcal{Y}})_S$ 与 $(\mathcal{D}_{\mathcal{Y}})_{S'}$ 的层论拉回等同起来。因此 $(\mathcal{D}_{\mathcal{Y}})_S$ 构成 $\mathcal{O}_{\mathcal{Y}}$-模 $\mathcal{D}_{\mathcal{Y}}$。实际上它是 1.1.5 意义下 $\mathcal{Y}$ 上的 $D$-模。而 $(D_{\mathcal{Y}})_S$ 构成 $\mathcal{Y}_{sm}$ 上结合代数的层 $D_{\mathcal{Y}}$，称为 $\mathcal{Y}$ 上\SourcePage{8}%
微分算子的层。层的嵌入 $D_{\mathcal{Y}}\hookrightarrow\mathcal{D}_{\mathcal{Y}}$ 给出整体截面空间之间的同构
\begin{equation}\tag{6}
\Gamma(\mathcal{Y},D_{\mathcal{Y}})\simeq\Gamma(\mathcal{Y},\mathcal{D}_{\mathcal{Y}}).
\end{equation}

\bdp{1.1.4} $\mathcal{O}_{\mathcal{Y}}$-模 $\mathcal{D}_{\mathcal{Y}}$ 带有按微分算子次数给出的自然滤链。$D_{\mathcal{Y}}$ 上的诱导滤链是代数滤链，使得 $\operatorname{gr}D_{\mathcal{Y}}$ 交换；因此 $\operatorname{gr}D_{\mathcal{Y}}$ 按通常方式成为泊松代数。

有显然的分次 $\mathcal{O}_{\mathcal{Y}}$-模满态射 $\operatorname{Sym}\Theta_{\mathcal{Y}}\to\operatorname{gr}D_{\mathcal{Y}}$。1.1.1 中的条件 (2) 保证它是同构。若此事实成立，则逆同构 $\operatorname{gr}D_{\mathcal{Y}}\simeq\operatorname{Sym}\Theta_{\mathcal{Y}}$ 诱导典范的泊松代数嵌入
\begin{equation}\tag{7}
\sigma_{\mathcal{Y}}:\operatorname{gr}D_{\mathcal{Y}}\hookrightarrow P_{\mathcal{Y}},
\end{equation}
称为象征映射。

\emph{注.} 在上面的叙述中，我们利用技术性条件 (2) 走了一条捷径。在 1.1.2--1.1.4 中真正应考虑的对象，是集中在次数 $\le 0$ 的复形（此时象征映射总有定义）；我们定义的朴素对象是它们的零次上同调。条件 (2) 蕴涵其余上同调为零，故不必为导出范畴费心（关于任意光滑叠 $\mathcal{Y}$ 的「真正」$\mathcal{D}_{\mathcal{Y}}$ 的定义，见 7.3.3）。

\bdp{1.1.5} $D$-模是光滑拓扑下的局部对象，因此光滑叠上 $D$-模的概念是清楚的\footnote{$D$-模的导出范畴的定义并不那么清楚；见 7.3。}。确切地说，对 $S\in\mathcal{Y}_{sm}$，$S$ 上左 $D$-模的范畴 $\mathcal{M}^{\ell}(S)$ 构成 $\mathcal{Y}_{sm}$ 上阿贝尔范畴的层 $\underline{\mathcal{M}}^{\ell}$（拉回函子是 $D$-模的通常拉回；由于 $\mathcal{Y}_{sm}$ 中的态射光滑，它们正合）。$\mathcal{Y}$ 上的 $D$-模是 $\underline{\mathcal{M}}^{\ell}$ 在 $\mathcal{Y}_{sm}$ 上的笛卡尔截面；它们构成阿贝尔范畴 $\mathcal{M}^{\ell}(\mathcal{Y})$。换言之，$\mathcal{Y}$ 上的 $D$-模是拟凝聚 $\mathcal{O}_{\mathcal{Y}}$-模 $M$，连同\SourcePage{9}%
每个 $\mathcal{O}_S$-模 $M_S$ 上相容的 $\mathcal{D}_S$-模结构（$S\in\mathcal{Y}_{sm}$）。通常的张量积使 $\mathcal{M}^{\ell}(\mathcal{Y})$ 成为张量范畴。凝聚、完整等 $D$-模可按显然方式定义。注意，$\mathcal{Y}$ 上的 $D$-模 $M$ 定义了 $\mathcal{Y}_{sm}$ 上结合代数的层 $\underline{\operatorname{End}}M$，且 $\underline{\operatorname{End}}M(S)=\underline{\operatorname{End}}M_S$。

例如，1.1.3 中我们定义了 $\mathcal{Y}$ 上的 $D$-模 $\mathcal{D}_{\mathcal{Y}}$；代数 $D_{\mathcal{Y}}$ 与 $\underline{\operatorname{End}}\mathcal{D}_{\mathcal{Y}}$ 反同构。

\bdp{1.1.6} 设 $\mathcal{L}$ 是 $\mathcal{Y}$ 上的线丛，$\lambda\in\mathbb{C}$。任意 $S\in\mathcal{Y}_{sm}$ 带有线丛 $\pi_S^{*}\mathcal{L}$。因此有 $\pi_S^{*}(\mathcal{L})^{\otimes\lambda}$-扭左 $D$-模的范畴 $\mathcal{M}^{\ell}(S)_{\mathcal{L}\lambda}$（例如见 [BB93]）。这些范畴构成 $\mathcal{Y}_{sm}$ 上阿贝尔范畴的层 $\underline{\mathcal{M}}^{\ell}_{\mathcal{L}\lambda}$。$\mathcal{Y}$ 上 $\mathcal{L}^{\otimes\lambda}$-扭 $D$-模的范畴 $\mathcal{M}^{\ell}(\mathcal{Y})_{\mathcal{L}\lambda}$，就是 $\underline{\mathcal{M}}^{\ell}_{\mathcal{L}\lambda}$ 的笛卡尔截面的范畴。有典范的完全忠实嵌入 $\mathcal{M}^{\ell}(\mathcal{Y})_{\mathcal{L}\lambda}\hookrightarrow\mathcal{M}^{\ell}(\mathcal{L}^{\cdot})$。它把 $\mathcal{Y}$ 上 $\mathcal{L}^{\otimes\lambda}$-扭 $D$-模与 $\mathcal{L}^{\cdot}$ 上 $\lambda$-单演 $D$-模等同起来。这里 $\mathcal{L}^{\cdot}$ 是对应于 $\mathcal{L}$ 的 $\mathbb{G}_m$-torsor（主齐性空间），即去掉零截面后的 $\mathcal{L}$。见 [BB93] 第 2 节。

我们留给读者去定义特出对象 $\mathcal{D}_{\mathcal{Y},\mathcal{L}\lambda}\in\mathcal{M}^{\ell}(\mathcal{Y})_{\mathcal{L}\lambda}$，以及 $\mathcal{Y}_{sm}$ 上滤结合代数的层 $D_{\mathcal{Y},\mathcal{L}\lambda}$。1.1.3--1.1.5 的全部事实都可无改动地搬到扭情形。

\bdp{1.1.7} 第 5 节将需要任意（不必光滑的）局部有限型代数叠上 $D$-模的概念。对概形而言此概念是熟知的（例如见 [Sa91]）。它对光滑拓扑是局部的，因此向叠的推广是直接的。

\bdsec{1.2. 若干熟知的构造}

\bdp{1.2.1} 设 $K$ 是作用在 $\mathbb{C}$ 上光滑概形 $S$ 的代数群。考虑商叠 $\mathcal{Y}=K\backslash S$。则 $S$ 是 $\mathcal{Y}$ 在 $\mathcal{Y}_{sm}$ 中的一个覆盖；而 $\mathcal{Y}$ 上的 $D$-模、线丛与扭 $D$-模，等同于 $S$ 上相应的 $K$-等变对象。$K$ 在 $T^{*}S$ 上的作用是哈密顿的，$T^{*}\mathcal{Y}$ 由 $T^{*}S$ 经哈密顿\SourcePage{10}%
约化得到，即 $T^{*}\mathcal{Y}=K\backslash\mu^{-1}(0)$，其中 $\mu:T^{*}S\to\mathfrak{k}^{*}$ 是矩映射，$\mathfrak{k}:=\operatorname{Lie}(K)$。$\Gamma(T^{*}\mathcal{Y},\mathcal{O}_{T^{*}\mathcal{Y}})$ 上的泊松结构，由把它与 $\Gamma(\mathcal{Y},P_{\mathcal{Y}})$ 等同而得到（见 1.1.2）；后者可用覆盖 $S\to\mathcal{Y}$ 来计算：
\begin{equation}\tag{8}
\Gamma(\mathcal{Y},P_{\mathcal{Y}})=\Gamma(S,\tilde{P}_S/I_S^{cl})^{\pi_0(K)}.
\end{equation}
这里 $\tilde{P}\subset\operatorname{Sym}\Theta_S$ 是理想 $I_S^{cl}:=(\operatorname{Sym}\Theta_S)\mathfrak{k}$ 的 $\{\,\}$-正规化子（$\mathfrak{k}$ 映到 $\Theta_S\subset\operatorname{Sym}\Theta_S$）。根据 1.1.3，
\begin{equation}\tag{9}
\Gamma(\mathcal{Y},D_{\mathcal{Y}})=\Gamma(S,\tilde{D}_S/I_S)^{\pi_0(K)}
\end{equation}
其中 $\tilde{D}_S\subset\mathcal{D}_S$ 是 $I_S:=\mathcal{D}_S\cdot\mathfrak{k}$ 的正规化子。

下面关于 $\mathcal{Y}$ 上象征、微分算子与 $D$-模的构造是有用的。

\bdp{1.2.2} 从 Harish-Chandra 对 $(\mathfrak{g},K)$ 出发。于是 $\mathfrak{g}$ 是李代数，带有 $K$ 的作用（称为伴随作用），并带有与 $K$ 的伴随作用相容的李代数嵌入 $\mathfrak{k}\hookrightarrow\mathfrak{g}$。设 $\tilde{P}_{(\mathfrak{g},K)}\subset\operatorname{Sym}\mathfrak{g}$ 是 $I^{cl}_{(\mathfrak{g},K)}:=(\operatorname{Sym}\mathfrak{g})\mathfrak{k}$ 的 $\{\,\}$-正规化子，$\tilde{D}_{(\mathfrak{g},K)}\subset U\mathfrak{g}$ 是 $I_{(\mathfrak{g},K)}:=(U\mathfrak{g})\mathfrak{k}$ 的正规化子。令
\begin{equation}\tag{10}
P_{(\mathfrak{g},K)}:=(\operatorname{Sym}(\mathfrak{g}/\mathfrak{k}))^{K}=(\tilde{P}_{(\mathfrak{g},K)}/I^{cl}_{(\mathfrak{g},K)})^{\pi_0(K)}
\end{equation}
\begin{equation}\tag{11}
D_{(\mathfrak{g},K)}:=(U\mathfrak{g}/(U\mathfrak{g})\mathfrak{k})^{K}=(\tilde{D}_{(\mathfrak{g},K)}/I_{(\mathfrak{g},K)})^{\pi_0(K)}.
\end{equation}
则 $P_{(\mathfrak{g},K)}$ 是泊松代数，$D_{(\mathfrak{g},K)}$ 是结合代数。$U\mathfrak{g}$ 上的标准滤链诱导 $D_{(\mathfrak{g},K)}$ 上的滤链，使 $\operatorname{gr}D_{(\mathfrak{g},K)}$ 交换。于是 $\operatorname{gr}D_{(\mathfrak{g},K)}$ 是泊松代数。有显然的泊松代数嵌入 $\sigma=\sigma_{(\mathfrak{g},K)}:\operatorname{gr}D_{(\mathfrak{g},K)}\hookrightarrow P_{(\mathfrak{g},K)}$。$(\mathfrak{g},K)$ 的局部量子化条件是说
\begin{equation}\tag{12}
\sigma_{(\mathfrak{g},K)}\ \text{是同构}.
\end{equation}

\SourcePage{11}

\emph{注.} 有时人们如下验证此条件。考虑显然的态射
\begin{equation}\tag{13}
a^{cl}:((\operatorname{Sym}\mathfrak{g})^{\mathfrak{g}})^{\pi_0(K)}\to P_{(\mathfrak{g},K)},\qquad
a:(\operatorname{Center}U\mathfrak{g})^{\pi_0(K)}\to D_{(\mathfrak{g},K)}.
\end{equation}
若 $a^{cl}$ 满，则 (12) 成立（因为 $\operatorname{gr}\,\operatorname{Center}U\mathfrak{g}=(\operatorname{Sym}\mathfrak{g})^{\mathfrak{g}}$）。事实上，若 $a^{cl}$ 满，则 $a$ 也满，故 $D_{(\mathfrak{g},K)}$ 交换。

\bdp{1.2.3} 现在假设处于 1.2.1 的情形，并且 $S$ 上的 $K$-作用延拓为 $(\mathfrak{g},K)$-作用。即有一个李代数态射 $\mathfrak{g}\to\Theta_{S'}$，它与 $S$ 上的 $K$-作用以显然的方式相容。比较 (8) 与 (10)、(9) 与 (11) 可见，态射 $\operatorname{Sym}\mathfrak{g}\to\operatorname{Sym}\Theta_S$ 与 $U\mathfrak{g}\to\mathcal{D}_S$ 诱导典范态射
% [Tutor Note: 原书 p.11 此处印作 $\Theta_{S'}$（$S$ 带撇），按上下文应为 $\Theta_S$；底本如此，保留原样。]
\begin{equation}\tag{14}
h^{cl}:P_{(\mathfrak{g},K)}\to\Gamma(\mathcal{Y},P_{\mathcal{Y}}),\qquad
h:D_{(\mathfrak{g},K)}\to\Gamma(\mathcal{Y},D_{\mathcal{Y}})
\end{equation}
它们分别是泊松代数与滤结合代数的态射。

若 $\mathcal{Y}$ 在 1.1.1 的意义下为良，则有象征映射 $\sigma_{\mathcal{Y}}:\operatorname{gr}D_{\mathcal{Y}}\hookrightarrow P_{\mathcal{Y}}$；上述态射是 $\sigma$-相容的：$h^{cl}\sigma_{(\mathfrak{g},K)}=\sigma_{\mathcal{Y}}\operatorname{gr}h$。我们数据的整体量子化条件是说
\begin{equation}\tag{15}
h\ \text{与滤链严格相容}.
\end{equation}
换言之，这意味着来自 $h(D_{(\mathfrak{g},K)})$ 的微分算子的象征落在 $h^{cl}\sigma_{(\mathfrak{g},K)}(\operatorname{gr}D_{(\mathfrak{g},K)})$ 中。若局部与整体量子化条件都满足，则微分算子代数 $h(D_{(\mathfrak{g},K)})$ 是象征代数 $h^{cl}(P_{(\mathfrak{g},K)})$ 的量子化：象征映射 $\sigma_{\mathcal{Y}}$ 给出同构 $\operatorname{gr}h(D_{(\mathfrak{g},K)})\simeq h^{cl}(P_{(\mathfrak{g},K)})$。

\emph{注.} 局部与整体量子化条件在某种意义下互补：局部条件说明 $D_{(\mathfrak{g},K)}$ 尽可能大，整体条件则意味着 $h(D_{(\mathfrak{g},K)})$ 尽可能小。

\SourcePage{12}

\bdp{1.2.4} 记 $\mathcal{M}(\mathfrak{g},K)$ 为 Harish-Chandra 模的范畴。有一对伴随函子（例如见 [BB93]）
\[
\Delta:\mathcal{M}(\mathfrak{g},K)\to\mathcal{M}^{\ell}(\mathcal{Y}),\qquad
\Gamma:\mathcal{M}^{\ell}(\mathcal{Y})\to\mathcal{M}(\mathfrak{g},K).
\]
确切地说，对 $\mathcal{Y}$ 上的 $D$-模 $M$，Harish-Chandra 模 $\Gamma(M)$ 是截面空间 $\Gamma(S,M_S)$，带有显然的 $(\mathfrak{g},K)$-作用（例如 $\mathfrak{g}$ 经由 $\mathfrak{g}\to\Theta_S\subset\mathcal{D}_S$ 作用）。对 $(\mathfrak{g},K)$-模 $V$，相应的 $K$-等变 $D$-模是 $\Delta(V)_S=\mathcal{D}_S\otimes_{U\mathfrak{g}}V$。

例如，考虑「真空」Harish-Chandra 模 $\operatorname{Vac}:=U\mathfrak{g}/(U\mathfrak{g})\mathfrak{k}$。对任意 $V\in\mathcal{M}(\mathfrak{g},K)$ 有 $\operatorname{Hom}(\operatorname{Vac},V)=V^{K}$。故有典范双射 $\operatorname{End}(\operatorname{Vac})\to\operatorname{Vac}^{K}=D_{(\mathfrak{g},K)}$（见 (11)），它实际上是代数的反同构。有显然的同构 $\Delta(\operatorname{Vac})=\mathcal{D}_{\mathcal{Y}}$，而映射 $\Delta:\operatorname{End}(\operatorname{Vac})\to\operatorname{End}(\mathcal{D}_{\mathcal{Y}})=\Gamma(\mathcal{Y},D_{\mathcal{Y}})^{\circ}$ 与 (14) 中的映射 $h$ 一致。

\bdp{1.2.5} 上述构造有扭版本。确切地说，假设把 $(\mathfrak{g},K)$ 用 $\mathbb{C}$ 作中心扩张，得到 $(\tilde{\mathfrak{g}},K)$，于是 $\mathbb{C}\subset\tilde{\mathfrak{g}}$，$\tilde{\mathfrak{g}}/\mathbb{C}=\mathfrak{g}$。记 $U'\mathfrak{g}$ 为 $U\tilde{\mathfrak{g}}$ 模去由中心元素 $\mathbf{1}-1$ 生成的理想所得的商，其中 $\mathbf{1}\in\mathbb{C}\subset\tilde{\mathfrak{g}}$。这是滤结合代数；把 $\operatorname{gr}U'\mathfrak{g}$ 与 $\operatorname{Sym}\mathfrak{g}$ 等同（作为泊松代数）。我们得到滤结合代数 $D'_{(\mathfrak{g},K)}:=(U'\mathfrak{g}/(U'\mathfrak{g})\mathfrak{k})^{K}$，带有嵌入 $\sigma:\operatorname{gr}D'_{(\mathfrak{g},K)}\hookrightarrow P_{(\mathfrak{g},K)}$。扭局部量子化条件说 $\sigma$ 是同构。注意 1.2.2 末的注在扭情形不成立，因为 $\operatorname{gr}\,\operatorname{Center}U'\mathfrak{g}$ 可能不等于 $(\operatorname{Sym}\mathfrak{g})^{\mathfrak{g}}$。

设 $\mathcal{L}$ 是 $S$ 上的线丛。假设 $S$ 上的 $(\mathfrak{g},K)$-作用提升为 $\mathcal{L}$ 上的 $(\tilde{\mathfrak{g}},K)$-作用，使得 $\mathbf{1}$ 乘以某个 $\lambda^{-1}$，其中 $\lambda\in\mathbb{C}^{*}$。等价地，有 $\mathcal{L}^{\cdot}$ 上的 $(\tilde{\mathfrak{g}},K)$-作用：它延拓 $K$-作用，与 $S$ 上的 $\mathfrak{g}$-作用相容，且 $\mathbf{1}$ 作用为 $-\lambda^{-1}t\partial_t\in\Theta_{\mathcal{L}^{\cdot}}$。令 $D'_{\mathcal{Y}}=\mathcal{D}_{\mathcal{Y},\mathcal{L}\lambda}$。有滤结合代数的态射 $h:D'_{(\mathfrak{g},K)}\to\Gamma(\mathcal{Y},D'_{\mathcal{Y}})$，使得 $\sigma\operatorname{gr}h=h^{cl}\sigma$。扭整体量子化条件说 $h$ 与滤链严格相容。

\SourcePage{13}

记 $\mathcal{M}(\mathfrak{g},K)'$ 为 $(\tilde{\mathfrak{g}},K)$-模范畴的完全子范畴，由那些 $\mathbf{1}$ 作用为恒等的 Harish-Chandra 模组成。在 $\mathcal{M}(\mathfrak{g},K)'$ 与 $\mathcal{M}^{\ell}(\mathcal{Y})_{\mathcal{L}\lambda}$ 之间有伴随函子 $\Delta,\Gamma$，其定义与非扭版本完全一样。再次，对 $\operatorname{Vac}':=U'\mathfrak{g}/(U'\mathfrak{g})\mathfrak{k}$ 有 $\Delta(\operatorname{Vac}')=\mathcal{D}_{\mathcal{Y},\mathcal{L}\lambda}$。代数 $\operatorname{End}(\operatorname{Vac}')$ 与 $D'_{(\mathfrak{g},K)}$ 反同构，且 $\Delta:\operatorname{End}(\operatorname{Vac}')\to\operatorname{End}\mathcal{D}_{\mathcal{Y},\mathcal{L}\lambda}=\Gamma(\mathcal{Y},D'_{\mathcal{Y}})$ 与 $h$ 一致。

\bdp{1.2.6} 无限维版本。设 $K$ 是 $\mathbb{C}$ 上的仿射群概形（于是 $K$ 是代数群的投射极限），它作用在概形 $S$ 上。假设以下条件：
\begin{equation}\tag{16}
\parbox{0.78\linewidth}{\raggedright
存在 $S$ 的 Zariski 开覆盖 $\{U_i\}$，使得每个 $U_i$ 都是 $K$-不变的；并且对某个正规群子概形 $K_i\subset K$（$K/K_i$ 为有限型），$U_i$ 是光滑概形 $T_i$ 上的主 $K_i$-丛（于是 $T_i=K_i\backslash U_i$）。}
\end{equation}
于是 fpqc-商 $\mathcal{Y}=K\backslash S$ 是光滑代数叠（它被开子叠 $(K/K_i)\backslash T_i$ 覆盖）。

我们来解释如何把 1.2.1--1.2.5 搬到这一情形。注意 $\mathfrak{k}=\operatorname{Lie}K$ 是有限维李代数的投射极限，因此是完全拓扑李代数。考虑层 $\Theta_S=\operatorname{Der}\mathcal{O}_S$，以及 Grothendieck 微分算子的层 $\mathcal{D}_S\subset\operatorname{End}_{\mathbb{C}}(\mathcal{O}_S)$。它们是完全拓扑李代数（相应地，结合代数）的层。确切地说，对仿射开集 $U\subset S$，$\Gamma(U,\Theta_S)$ 与 $\Gamma(U,\mathcal{D}_S)$ 中开子空间的基本邻域，由 $\Gamma(U,\mathcal{O}_U)$ 的有限生成子代数的零化子组成。$\Theta_S$ 上的拓扑定义了 $\operatorname{Sym}\Theta_S$ 上的拓扑；记 $\overline{\operatorname{Sym}}\Theta_S$ 为完备化代数。这是拓扑泊松代数的层。设 $I_S^{cl}\subset\overline{\operatorname{Sym}}\Theta_S$ 是理想 $(\overline{\operatorname{Sym}}\Theta_S)\mathfrak{k}$ 的闭包，$\tilde{P}_S\subset\overline{\operatorname{Sym}}\Theta_S$ 是它的 $\{\,\}$-正规化子。类似地，设 $I_S\subset\mathcal{D}_S$ 是理想 $\mathcal{D}_S\cdot\mathfrak{k}$ 的闭包，$\tilde{D}_S$ 是它的正规化子。则 (8)、(9) 中的公式仍然成立。

\SourcePage{14}

在 Harish-Chandra 对 $(\mathfrak{g},K)$ 的定义中，我们假设：对任意 $\operatorname{Ad}(K)$-不变开子空间 $\mathfrak{a}\subset k$，$K$ 在 $\mathfrak{g}/\mathfrak{a}$ 上的作用是代数的。
% [Tutor Note: 原书 p.14 此处印作斜体 $k$，按上下文应为 $\mathfrak{k}$；底本如此，保留原样。]

于是 $\mathfrak{g}$ 是完全拓扑李代数（$\mathfrak{g}$ 上的拓扑使得 $\mathfrak{k}\subset\mathfrak{g}$ 是开嵌入）。代数 $\operatorname{Sym}\mathfrak{g}$、$U\mathfrak{g}$ 带有自然拓扑，由开理想 $(\operatorname{Sym}\mathfrak{g})\mathfrak{a}$、$(U\mathfrak{g})\mathfrak{a}$ 定义；其中 $\mathfrak{a}\subset\mathfrak{g}$ 是开子代数。记相应的完备化为 $\overline{\operatorname{Sym}}\mathfrak{g}$、$\overline{U}\mathfrak{g}$。设 $I^{cl}_{(\mathfrak{g},K)}\subset\overline{\operatorname{Sym}}\mathfrak{g}$ 是理想 $(\overline{\operatorname{Sym}}\mathfrak{g})\mathfrak{k}$ 的闭包，$\tilde{P}_{(\mathfrak{g},K)}$ 是它的 $\{\,\}$-正规化子）。类似地，有 $I_{(\mathfrak{g},K)}\subset\tilde{D}_{(\mathfrak{g},K)}\subset\overline{U}\mathfrak{g}$。现在用公式 (10)、(11) 定义 $P_{(\mathfrak{g},K)}$、$D_{(\mathfrak{g},K)}$。1.2.2--1.2.5 的其余部分仍然成立，但 1.2.2 末的注除外。该注应如下修改。
% [Tutor Note: 原书 p.14「be its $\{\}$-normalizer).」一句多出一个右括号；底本如此，保留原样。]

\bdp{1.2.7} 代数 $\overline{\operatorname{Sym}}\mathfrak{g}$ 与 $\overline{U}\mathfrak{g}$ 带有通常的环滤链 $\overline{\operatorname{Sym}}_n\mathfrak{g}=\bigoplus_{0\le i\le n}\overline{\operatorname{Sym}}^{i}\mathfrak{g}$ 与 $\overline{U}_i\mathfrak{g}$。然而在无限维情形，这些滤链各项的并集并不与整个代数重合。有通常的同构 $\bar{\sigma}_{\mathfrak{g}}:\operatorname{gr}_i\overline{U}\mathfrak{g}\simeq\overline{\operatorname{Sym}}^{i}\mathfrak{g}$。对 $\overline{\operatorname{Sym}}\Theta_S$ 与 $\mathcal{D}_S$ 也有同样的事实。

1.2.2 末的态射 $a^{cl}$、$a$ 按显然方式延拓为态射
\begin{equation}\tag{17}
\bar{a}^{cl}:((\overline{\operatorname{Sym}}\mathfrak{g})^{\mathfrak{g}})^{\pi_0(K)}\to P_{(\mathfrak{g},K)},\qquad
\bar{a}:(\operatorname{Center}\overline{U}\mathfrak{g})^{\pi_0(K)}\to D_{(\mathfrak{g},K)}.
\end{equation}
1.2.2 中的局部量子化条件 (12) 与 $\bar{a}$ 的满性，都可由 $\bar{a}^{cl}\bar{\sigma}_{\mathfrak{g}}:\operatorname{gr}(\operatorname{Center}\overline{U}\mathfrak{g})^{\pi_0(K)}\to P_{(\mathfrak{g},K)}$ 的满性推出。扭情形也同样。注意等式 $\operatorname{gr}\,\operatorname{Center}\overline{U}\mathfrak{g}=(\operatorname{Sym}\mathfrak{g})^{\mathfrak{g}}$ 未必成立（即使在非扭情形！）。
